% ============================================================================
%  PRECÁLCULO — Sesión 1
%  Semana 1: 10 al 15 de agosto de 2026
%  Presentación del curso · 1.1 Conjuntos numéricos y sus propiedades
%  Duración prevista: 3 horas
% ============================================================================
\documentclass[aspectratio=169,10pt]{beamer}

\usepackage[utf8]{inputenc}
\IfFileExists{spanish.ldf}{\usepackage[spanish,es-noshorthands]{babel}}{}
\usepackage{amsmath,amssymb}
\usepackage{booktabs}
\usepackage{tikz}

\usetheme{Uniminuto}

\uniasignatura{Precálculo}
\uniprofesor{Andrés Felipe Puerta Velez}
\unicorreo{andres.puerta-v@uniminuto.edu.co}
\uninrc{MATE1161}
\unisemana{Sesión 1 — Semana 1 · 10 al 15 de agosto de 2026}

\newcommand{\listatiny}{\setbeamerfont{itemize/enumerate body}{size=\tiny}\tiny}

\begin{document}

\portada

\begin{frame}{Orden del día}
  \begin{enumerate}
    \item Llamado a lista, presentación del grupo y elección del representante
    \item ¿Para qué sirve Precálculo en Ingeniería de Software?
    \item Temario, cronograma, evaluación y acuerdos mínimos
    \item \azul{Tema 1.1: los conjuntos numéricos y sus propiedades}
    \item Orden en $\mathbb{R}$: desigualdades, intervalos y valor absoluto
    \item Ejercicios guiados, taller y cierre
  \end{enumerate}

  \vspace{0.25cm}
  {\footnotesize\color{uniGris} Sesión de 3 horas. Traiga cuaderno cuadriculado
  y lápiz: \destacar{en este curso se escribe a mano}, porque la única forma de
  aprender a operar es operando.}
\end{frame}

\begin{frame}{¿Para qué sirve Precálculo en Ingeniería de Software?}
  \scriptsize
  \begin{center}
  \begin{tabular}{@{}p{5.0cm}p{6.8cm}@{}}
    \toprule
    \textbf{Lo que verá más adelante} & \textbf{Lo que exige de este curso} \\
    \midrule
    Cálculo diferencial e integral &
      Álgebra fluida: factorizar, simplificar, despejar sin equivocarse \\
    \addlinespace
    Estructuras discretas y lógica &
      Conjuntos, operaciones y notación formal \\
    \addlinespace
    Complejidad de algoritmos &
      Potencias y \azul{logaritmos}: $O(n\log n)$ no se entiende sin ellos \\
    \addlinespace
    Gráficos, videojuegos y visión &
      Trigonometría: senos, cosenos y ángulos \\
    \addlinespace
    Criptografía y estructuras de datos &
      Aritmética entera, divisibilidad, potencias de 2 \\
    \bottomrule
  \end{tabular}
  \end{center}

  \vspace{0.1cm}
  \begin{alertblock}{La razón de fondo}
    \scriptsize
    Un programa es una \destacar{secuencia de operaciones exactas}. Quien no
    opera bien a mano, no detecta cuándo su código está calculando mal. Este
    curso es el que le da esa detección automática.
  \end{alertblock}
\end{frame}

\begin{frame}{El curso en tres unidades}
  \scriptsize
  \begin{center}
  \begin{tabular}{@{}p{1.2cm}p{4.4cm}p{6.2cm}@{}}
    \toprule
    & \textbf{Unidad} & \textbf{Temas} \\
    \midrule
    1 & Campos numéricos &
      Conjuntos numéricos; enteros; racionales; conjuntos; razones y
      proporciones \\
    \addlinespace
    2 & Potenciación, radicación y logaritmación &
      Potencias; radicación; racionalización; logaritmos y sus propiedades \\
    \addlinespace
    3 & Álgebra y trigonometría &
      Operaciones algebraicas; productos notables; factorización; fracciones
      algebraicas; funciones trigonométricas \\
    \bottomrule
  \end{tabular}
  \end{center}

  \vspace{0.15cm}
  \begin{alertblock}{Cómo se conectan}
    \scriptsize
    La unidad 1 fija \destacar{con qué números} se trabaja; la unidad 2, las
    \destacar{operaciones} que se hacen con ellos; la unidad 3, cómo operar
    cuando en lugar de números hay \destacar{letras}. Cada una se apoya
    entera en la anterior: \azul{aquí no hay temas sueltos}.
  \end{alertblock}
\end{frame}

\begin{frame}{Evaluación: tres exámenes, nada más}
  \footnotesize
  \begin{center}
  \begin{tabular}{@{}llr@{}}
    \toprule
    \textbf{Examen} & \textbf{Fecha} & \textbf{Peso} \\
    \midrule
    Examen 1 — Unidad 1 & Semana 6: 14 al 19 de septiembre & 35\% \\
    Examen 2 — Unidad 2 & Semana 12: 26 al 31 de octubre   & 35\% \\
    Examen 3 — Unidad 3 & Semana 16: 23 al 28 de noviembre & 30\% \\
    \midrule
    \textbf{Total} & & \textbf{100\%} \\
    \bottomrule
  \end{tabular}
  \end{center}

  \vspace{0.1cm}
  \begin{block}{Cómo se presentan}
    \tiny
    Individuales, en \azul{Aulas Virtuales (Moodle)}, de forma
    \azul{presencial en el aula} y con tiempo límite. Cada examen evalúa la
    \destacar{totalidad} de los temas de su unidad.
  \end{block}

  \vspace{0.08cm}
  \begin{alertblock}{Lo que esto implica}
    \tiny
    No hay quices ni talleres calificables que suban la nota. \destacar{La
    práctica semanal no se califica, pero es lo único que hace pasar el
    examen.}
  \end{alertblock}
\end{frame}

\begin{frame}{Acuerdos mínimos del curso}
  \scriptsize
  \begin{columns}[T,onlytextwidth]
    \begin{column}{0.48\textwidth}
      \begin{block}{De trabajo}
        \tiny
        \begin{itemize}\setlength{\itemsep}{-0.04cm}
          \item Se resuelve \azul{a mano}, mostrando el procedimiento.
          \item Cada sesión abre con las dudas de los ejercicios propuestos.
          \item La calculadora se usa para verificar, no para reemplazar el
                procedimiento.
          \item Preguntar en clase es parte del trabajo, no una interrupción.
        \end{itemize}
      \end{block}
    \end{column}
    \begin{column}{0.48\textwidth}
      \begin{alertblock}{De convivencia}
        \tiny
        \begin{itemize}\setlength{\itemsep}{-0.04cm}
          \item Puntualidad al inicio y después del descanso.
          \item Celular en silencio; se usa solo cuando el ejercicio lo pida.
          \item Respeto por el error ajeno: \destacar{equivocarse en el
                tablero es la forma más rápida de aprender}.
        \end{itemize}
      \end{alertblock}
    \end{column}
  \end{columns}

  \vspace{0.1cm}
  \begin{exampleblock}{El compromiso que de verdad importa}
    \tiny
    \destacar{Media hora diaria de ejercicios} rinde más que cinco horas la
    noche anterior al examen. Las matemáticas no se entienden mirando: se
    entienden con el lápiz en la mano.
  \end{exampleblock}
\end{frame}

\seccion{Unidad 1. Campos numéricos}

\begin{frame}{1.1 Los conjuntos numéricos}
  \centering
  \begin{tikzpicture}[scale=0.82]
    \draw[uniGris,fill=uniGris!8,rounded corners] (0,0) rectangle (10.4,4.1);
    \node[font=\scriptsize\bfseries,color=uniGris] at (9.6,3.8) {$\mathbb{R}$};
    \draw[uniAzul,fill=uniAzul!10,rounded corners] (0.25,0.25) rectangle (6.9,3.5);
    \node[font=\scriptsize\bfseries,color=uniAzul] at (6.4,3.2) {$\mathbb{Q}$};
    \draw[uniNaranja,fill=uniNaranja!12,rounded corners] (0.5,0.5) rectangle (4.9,2.9);
    \node[font=\scriptsize\bfseries,color=uniNaranja] at (4.4,2.6) {$\mathbb{Z}$};
    \draw[uniTitulo,fill=uniTitulo!12,rounded corners] (0.75,0.75) rectangle (3.2,2.3);
    \node[font=\scriptsize\bfseries,color=uniTitulo] at (2.8,2.0) {$\mathbb{N}$};
    \node[font=\tiny] at (1.9,1.4) {$0,\,1,\,2,\,3,\dots$};
    \node[font=\tiny] at (4.05,1.4) {$-1,\,-2,\dots$};
    \node[font=\tiny,align=center] at (5.9,1.7) {$\tfrac{1}{2},\ -\tfrac{7}{3}$\\ $0{,}25;\ 0{,}\overline{3}$};
    \draw[uniAmarillo!70!black,fill=uniAmarillo!18,rounded corners] (7.05,0.25) rectangle (10.15,3.5);
    \node[font=\scriptsize\bfseries,color=uniAmarillo!70!black] at (9.7,3.2) {$\mathbb{I}$};
    \node[font=\tiny,align=center,color=uniGris] at (8.6,1.875)
      {(irracionales)\\[0.15em] $\sqrt{2},\ \pi,\ e,\ \varphi$};
  \end{tikzpicture}

  \vspace{0.1cm}
  \begin{alertblock}{La cadena de inclusiones}
    \tiny
    \[ \mathbb{N} \subset \mathbb{Z} \subset \mathbb{Q} \subset \mathbb{R}
       \qquad\text{y}\qquad \mathbb{R} = \mathbb{Q} \cup \mathbb{I},
       \qquad \mathbb{Q} \cap \mathbb{I} = \varnothing \]
    Cada conjunto \destacar{contiene} al anterior: todo natural es entero, todo
    entero es racional y todo racional es real. \azul{Lo contrario no es
    cierto.}
  \end{alertblock}
\end{frame}

\begin{frame}{Naturales y enteros: por qué hubo que ampliarlos}
  \scriptsize
  \begin{center}
  \begin{tabular}{@{}p{1.6cm}p{3.6cm}p{3.0cm}p{3.4cm}@{}}
    \toprule
    \textbf{Conjunto} & \textbf{Qué contiene} & \textbf{Qué se puede} &
    \textbf{Qué falla} \\
    \midrule
    $\mathbb{N}$ & $0, 1, 2, 3, \dots$ &
      Sumar y multiplicar & $3 - 5$ no existe \\
    \addlinespace
    $\mathbb{Z}$ & $\dots, -2, -1, 0, 1, 2, \dots$ &
      Sumar, restar y multiplicar & $3 \div 5$ no existe \\
    \addlinespace
    $\mathbb{Q}$ & $\dfrac{a}{b}$ con $a,b \in \mathbb{Z}$, $b \ne 0$ &
      Las cuatro operaciones & $\sqrt{2}$ no existe \\
    \addlinespace
    $\mathbb{R}$ & Todos los puntos de la recta &
      Además, raíces de positivos & $\sqrt{-1}$ no existe \\
    \bottomrule
  \end{tabular}
  \end{center}

  \vspace{0.1cm}
  \begin{alertblock}{La lógica de toda la historia}
    \scriptsize
    Cada conjunto nuevo aparece porque una operación \destacar{se salía} del
    anterior. Es el mismo patrón que en programación: se amplía el tipo de dato
    cuando el que se tiene \azul{ya no alcanza}. Y en la siguiente asignatura
    aparecerán los complejos $\mathbb{C}$, justamente para que $\sqrt{-1}$
    tenga sentido.
  \end{alertblock}
\end{frame}

\begin{frame}{Racionales: la forma de reconocerlos}
  \scriptsize
  \begin{columns}[T,onlytextwidth]
    \begin{column}{0.46\textwidth}
      \begin{block}{Definición}
        \scriptsize
        \[ \mathbb{Q} = \left\{ \frac{a}{b} \;:\; a, b \in \mathbb{Z},\
           b \neq 0 \right\} \]
        Todo racional se puede escribir como \azul{cociente de dos enteros}.
      \end{block}
    \end{column}
    \begin{column}{0.50\textwidth}
      \begin{exampleblock}{Su expresión decimal}
        \scriptsize
        Un número es racional \azul{si y solo si} su decimal es
        \destacar{finito} o \destacar{periódico}:
        \[ \tfrac{1}{4} = 0{,}25 \qquad \tfrac{1}{3} = 0{,}\overline{3}
           \qquad \tfrac{5}{11} = 0{,}\overline{45} \]
      \end{exampleblock}
    \end{column}
  \end{columns}

  \vspace{0.08cm}
  \begin{block}{De decimal periódico a fracción: el truco que hay que saber}
    \tiny
    Sea $x = 0{,}\overline{45}$. Como el periodo tiene 2 cifras, se multiplica
    por $100$:
    \[ 100x = 45{,}\overline{45} \qquad x = 0{,}\overline{45} \]
    Restando: \ $99x = 45 \Rightarrow x = \dfrac{45}{99} = \destacar{\dfrac{5}{11}}$
  \end{block}

  {\tiny\color{uniGris} Regla general: se multiplica por $10^{k}$, con $k$ el
  número de cifras del periodo, y se resta la ecuación original.}
\end{frame}

\begin{frame}{Irracionales: los que no se pueden escribir como fracción}
  \scriptsize
  \begin{columns}[T,onlytextwidth]
    \begin{column}{0.44\textwidth}
      \begin{block}{Los más conocidos}
        \tiny
        \[ \sqrt{2} = 1{,}41421356\dots \]
        \[ \pi = 3{,}14159265\dots \]
        \[ e = 2{,}71828182\dots \]
        \[ \varphi = \tfrac{1+\sqrt{5}}{2} = 1{,}61803398\dots \]
        Decimal \destacar{infinito y no periódico}.
      \end{block}
    \end{column}
    \begin{column}{0.52\textwidth}
      \begin{exampleblock}{Por qué $\sqrt{2}$ no es racional}
        \tiny
        Suponga que sí: $\sqrt{2} = \frac{a}{b}$ ya simplificada. Al elevar
        al cuadrado:
        \[ 2 = \frac{a^2}{b^2} \;\Rightarrow\; a^2 = 2b^2 \]
        Entonces $a^2$ es par, luego $a$ es par: $a = 2k$. Sustituyendo:
        \[ 4k^2 = 2b^2 \;\Rightarrow\; b^2 = 2k^2 \]
        Entonces $b$ también es par. Pero si ambos son pares,
        \destacar{la fracción no estaba simplificada}: contradicción.
      \end{exampleblock}
    \end{column}
  \end{columns}

  \vspace{0.02cm}
  \begin{alertblock}{Dos advertencias}
    \tiny
    (1) $3{,}1416$ \destacar{no es} $\pi$: es una aproximación racional. \quad
    (2) No toda raíz es irracional: $\sqrt{9} = 3$ es natural. Solo lo es
    cuando el radicando \azul{no es un cuadrado perfecto}.
  \end{alertblock}
\end{frame}

\begin{frame}{Ejercicio 1 — en clase}
  \begin{alertblock}{Individual (10 minutos)}
    \scriptsize
    Para cada número, marque \destacar{todos} los conjuntos a los que
    pertenece: $\mathbb{N}$, $\mathbb{Z}$, $\mathbb{Q}$, $\mathbb{I}$,
    $\mathbb{R}$.

    \vspace{0.25cm}
    \centering
    \begin{tabular}{@{}llll@{}}
      1. \ $7$ & 2. \ $-4$ & 3. \ $\dfrac{3}{5}$ & 4. \ $\sqrt{16}$ \\[0.6em]
      5. \ $\sqrt{7}$ & 6. \ $0$ & 7. \ $-\dfrac{9}{3}$ & 8. \ $0{,}75$ \\[0.6em]
      9. \ $0{,}\overline{6}$ & 10. \ $\pi$ & 11. \ $\dfrac{22}{7}$ &
      12. \ $-\sqrt{25}$ \\
    \end{tabular}

    \vspace{0.25cm}
    \raggedright
    Al terminar: \destacar{¿cuál de los doce es el más engañoso y por qué?}
  \end{alertblock}
\end{frame}

\begin{frame}{Ejercicio 1 — solución}
  \tiny
  \begin{center}
  \begin{tabular}{@{}clcccc@{}}
    \toprule
    \textbf{\#} & \textbf{Número} & $\mathbb{N}$ & $\mathbb{Z}$ &
    $\mathbb{Q}$ & $\mathbb{I}$ \\
    \midrule
     1 & $7$                          & Sí & Sí & Sí & No \\
     2 & $-4$                         & No & Sí & Sí & No \\
     3 & $3/5$                        & No & No & Sí & No \\
     4 & $\sqrt{16} = 4$              & Sí & Sí & Sí & No \\
     5 & $\sqrt{7} = 2{,}6457\dots$   & No & No & No & Sí \\
     6 & $0$                          & Sí & Sí & Sí & No \\
     7 & $-9/3 = -3$                  & No & Sí & Sí & No \\
     8 & $0{,}75 = 3/4$               & No & No & Sí & No \\
     9 & $0{,}\overline{6} = 2/3$     & No & No & Sí & No \\
    10 & $\pi$                        & No & No & No & Sí \\
    11 & $22/7 = 3{,}\overline{142857}$ & No & No & Sí & No \\
    12 & $-\sqrt{25} = -5$            & No & Sí & Sí & No \\
    \bottomrule
  \end{tabular}
  \end{center}
  {\tiny\color{uniGris} Todos son reales: la columna $\mathbb{R}$ sería ``Sí''
  en las doce filas.}

  \vspace{0.05cm}
  \begin{alertblock}{El más engañoso: el 11}
    \tiny
    $22/7$ se usa como aproximación de $\pi$, pero \destacar{es racional}: es
    un cociente de enteros y su decimal es periódico. \azul{Parecerse a un
    irracional no vuelve irracional a un número.} Los otros dos tramposos son
    el 4 y el 12: llevan raíz, pero el radicando es cuadrado perfecto.
  \end{alertblock}
\end{frame}

\begin{frame}{Propiedades de las operaciones en $\mathbb{R}$}
  \tiny
  \begin{center}
  \begin{tabular}{@{}p{2.3cm}p{4.2cm}p{4.6cm}@{}}
    \toprule
    \textbf{Propiedad} & \textbf{Suma} & \textbf{Multiplicación} \\
    \midrule
    Clausurativa & $a+b \in \mathbb{R}$ & $a\cdot b \in \mathbb{R}$ \\
    Conmutativa  & $a+b = b+a$ & $a\cdot b = b\cdot a$ \\
    Asociativa   & $(a+b)+c = a+(b+c)$ & $(a\cdot b)\cdot c = a\cdot(b\cdot c)$ \\
    Modulativa   & $a+0 = a$ \quad (neutro: $0$) &
      $a\cdot 1 = a$ \quad (neutro: $1$) \\
    Invertiva    & $a+(-a) = 0$ &
      $a\cdot \dfrac{1}{a} = 1$ \quad ($a \neq 0$) \\
    \addlinespace
    Distributiva & \multicolumn{2}{l}{$a\cdot(b+c) = a\cdot b + a\cdot c$} \\
    \bottomrule
  \end{tabular}
  \end{center}

  \vspace{0.1cm}
  \begin{alertblock}{Las dos excepciones que hay que recordar}
    \tiny
    (1) La \destacar{resta y la división no son conmutativas ni asociativas}:
    $5-3 \neq 3-5$ y $(8\div4)\div2 \neq 8\div(4\div2)$. \quad
    (2) El \destacar{cero no tiene inverso multiplicativo}: por eso dividir
    entre cero no está definido, ni en matemáticas ni en su código.
  \end{alertblock}

  \vspace{0.05cm}
  {\tiny\color{uniGris} La distributiva es la que más se usa en todo el
  semestre: es la que permite ``abrir'' paréntesis y, leída al revés, la que
  permite \azul{factorizar} en la unidad 3.}
\end{frame}

\begin{frame}{Ejercicio 2 — en clase}
  \begin{alertblock}{En parejas (10 minutos)}
    \scriptsize
    Indique qué propiedad justifica cada igualdad:

    \vspace{0.25cm}
    \begin{columns}[T,onlytextwidth]
      \begin{column}{0.48\textwidth}
        \begin{enumerate}\setlength{\itemsep}{0.05cm}
          \item $7 + 12 = 12 + 7$
          \item $3(x+5) = 3x + 15$
          \item $(2+9)+8 = 2+(9+8)$
          \item $-6 + 6 = 0$
          \item $\sqrt{5}\cdot 1 = \sqrt{5}$
        \end{enumerate}
      \end{column}
      \begin{column}{0.48\textwidth}
        \begin{enumerate}\setlength{\itemsep}{0.05cm}
          \setcounter{enumi}{5}
          \item $\tfrac{4}{7}\cdot\tfrac{7}{4} = 1$
          \item $(5\cdot 3)\cdot 2 = 5\cdot(3\cdot 2)$
          \item $8 + 0 = 8$
          \item $2a + 2b = 2(a+b)$
          \item $\pi \cdot e = e \cdot \pi$
        \end{enumerate}
      \end{column}
    \end{columns}
  \end{alertblock}
\end{frame}

\begin{frame}{Ejercicio 2 — solución}
  \scriptsize
  \begin{columns}[T,onlytextwidth]
    \begin{column}{0.48\textwidth}
      \begin{block}{Respuestas 1 a 5}
        \scriptsize
        \begin{enumerate}\setlength{\itemsep}{-0.02cm}
          \item Conmutativa de la suma
          \item Distributiva
          \item Asociativa de la suma
          \item Invertiva de la suma (opuesto)
          \item Modulativa del producto
        \end{enumerate}
      \end{block}
    \end{column}
    \begin{column}{0.48\textwidth}
      \begin{block}{Respuestas 6 a 10}
        \scriptsize
        \begin{enumerate}\setlength{\itemsep}{-0.02cm}
          \setcounter{enumi}{5}
          \item Invertiva del producto (recíproco)
          \item Asociativa del producto
          \item Modulativa de la suma
          \item Distributiva \azul{leída al revés}
          \item Conmutativa del producto
        \end{enumerate}
      \end{block}
    \end{column}
  \end{columns}

  \vspace{0.1cm}
  \begin{alertblock}{El punto 9 es el importante}
    \scriptsize
    $2a + 2b = 2(a+b)$ es la distributiva usada de derecha a izquierda: eso
    \destacar{es factorizar}. Cuando lleguemos a la unidad 3 no será una
    técnica nueva, sino esta misma propiedad leída al revés.
  \end{alertblock}
\end{frame}

\begin{frame}{Orden en $\mathbb{R}$: la recta numérica}
  \centering
  \begin{tikzpicture}[scale=1.0]
    \draw[uniGris,->] (-0.4,0) -- (11.6,0);
    \foreach \x/\l in {0.5/-3, 2/-2, 3.5/-1, 5/0, 6.5/1, 8/2, 9.5/3, 11/4}
      \draw[uniGris] (\x,-0.12) -- (\x,0.12)
        node[below=3pt,font=\tiny,color=uniGris] {$\l$};
    \fill[uniAzul] (7.1,0) circle (2.2pt)
      node[above=3pt,font=\tiny,color=uniAzul] {$\sqrt{2}$};
    \fill[uniTitulo] (9.7,0) circle (2.2pt)
      node[above=3pt,font=\tiny,color=uniTitulo] {$\pi$};
    \fill[uniNaranja] (4.25,0) circle (2.2pt)
      node[above=3pt,font=\tiny,color=uniNaranja] {$-\tfrac{1}{2}$};
  \end{tikzpicture}

  \vspace{0.3cm}
  \begin{alertblock}{Lo que dice la recta}
    \scriptsize
    A cada punto le corresponde \destacar{un} número real y a cada número real
    \destacar{un} punto: no quedan huecos. Y el orden es visual:
    \azul{$a < b$ significa que $a$ está a la izquierda de $b$}. Por eso
    $-5 < -2$, aunque $5 > 2$.
  \end{alertblock}

  \vspace{0.1cm}
  \begin{exampleblock}{Reglas de las desigualdades}
    \scriptsize
    Si $a < b$: \quad $a + c < b + c$ (sumar no cambia nada) \quad y \quad
    $ac < bc$ si $c > 0$, \ pero \ \destacar{$ac > bc$ si $c < 0$}.
    \vspace{0.05cm}
    {\tiny\color{uniGris} Multiplicar por un negativo \azul{invierte} la
    desigualdad: es el error número uno al resolver inecuaciones.}
  \end{exampleblock}
\end{frame}

\begin{frame}{Intervalos y valor absoluto}
  \scriptsize
  \begin{center}
  \begin{tabular}{@{}p{2.2cm}p{3.0cm}p{3.0cm}p{3.4cm}@{}}
    \toprule
    \textbf{Nombre} & \textbf{Notación} & \textbf{Desigualdad} &
    \textbf{Se lee} \\
    \midrule
    Abierto & $(a,\,b)$ & $a < x < b$ & Sin los extremos \\
    Cerrado & $[a,\,b]$ & $a \le x \le b$ & Con los dos extremos \\
    Semiabierto & $[a,\,b)$ & $a \le x < b$ & Con $a$, sin $b$ \\
    Infinito & $[a,\,\infty)$ & $x \ge a$ & $\infty$ \azul{nunca} se cierra \\
    \bottomrule
  \end{tabular}
  \end{center}

  \vspace{0.1cm}
  \begin{columns}[T,onlytextwidth]
    \begin{column}{0.46\textwidth}
      \begin{block}{Valor absoluto}
        \tiny
        \[ |x| = \begin{cases} x & \text{si } x \ge 0 \\
                              -x & \text{si } x < 0 \end{cases} \]
        Es la \azul{distancia al cero}, y por eso nunca es negativo:
        $|{-7}| = 7$, \ $|7| = 7$.
      \end{block}
    \end{column}
    \begin{column}{0.50\textwidth}
      \begin{alertblock}{El error clásico}
        \tiny
        $|-x| \neq -x$ en general. Si $x = -3$, entonces $-x = 3$ y
        $|-x| = |3| = 3$. \destacar{El signo menos delante de la letra no
        significa ``número negativo''}: significa ``el opuesto de''.
      \end{alertblock}
    \end{column}
  \end{columns}
\end{frame}

\begin{frame}{Ejercicio 3 — en clase}
  \begin{alertblock}{Individual (12 minutos)}
    \scriptsize
    \textbf{Parte A.} Escriba cada conjunto en notación de intervalo y
    represéntelo en la recta:

    \vspace{0.15cm}
    \begin{columns}[T,onlytextwidth]
      \begin{column}{0.48\textwidth}
        \begin{enumerate}\setlength{\itemsep}{-0.02cm}
          \item $-2 \le x < 5$
          \item $x > 3$
          \item $-1 < x \le 4$
        \end{enumerate}
      \end{column}
      \begin{column}{0.48\textwidth}
        \begin{enumerate}\setlength{\itemsep}{-0.02cm}
          \setcounter{enumi}{3}
          \item $x \le 0$
          \item Todos los reales
        \end{enumerate}
      \end{column}
    \end{columns}

    \vspace{0.2cm}
    \textbf{Parte B.} Calcule:
    \[ |{-9}| \qquad |4 - 11| \qquad |3| - |{-3}| \qquad
       |{-2}| \cdot |{-5}| \qquad |x| \text{ si } x = -\tfrac{2}{3} \]
  \end{alertblock}
\end{frame}

\begin{frame}{Ejercicio 3 — solución}
  \scriptsize
  \begin{columns}[T,onlytextwidth]
    \begin{column}{0.44\textwidth}
      \begin{block}{Parte A}
        \tiny
        \begin{enumerate}\setlength{\itemsep}{-0.02cm}
          \item $[-2,\,5)$
          \item $(3,\,\infty)$
          \item $(-1,\,4]$
          \item $(-\infty,\,0]$
          \item $(-\infty,\,\infty) = \mathbb{R}$
        \end{enumerate}
      \end{block}
    \end{column}
    \begin{column}{0.52\textwidth}
      \begin{block}{Parte B}
        \tiny
        \[ |{-9}| = 9 \qquad |4-11| = |{-7}| = 7 \]
        \[ |3| - |{-3}| = 3 - 3 = 0 \]
        \[ |{-2}|\cdot|{-5}| = 2 \cdot 5 = 10 \]
        \[ \left|-\tfrac{2}{3}\right| = \tfrac{2}{3} \]
      \end{block}
    \end{column}
  \end{columns}

  \vspace{0.08cm}
  \begin{alertblock}{Dos trampas frecuentes}
    \tiny
    En $|4-11|$ \destacar{primero se opera adentro} y después se toma el valor
    absoluto: no es $|4| - |11| = -7$. \quad
    Y en el punto 5 de la parte A, los infinitos \destacar{siempre} llevan
    paréntesis: $[-\infty, \infty]$ está mal escrito, porque $\infty$ no es un
    número que se pueda alcanzar.
  \end{alertblock}
\end{frame}

\begin{frame}{Ejercicio 4 — en clase (taller principal)}
  \begin{alertblock}{En grupos de 3 (20 minutos)}
    \scriptsize
    \begin{enumerate}\setlength{\itemsep}{0.02cm}
      \item Convierta a fracción: \ $0{,}\overline{7}$, \quad
            $0{,}\overline{27}$, \quad $1{,}\overline{3}$.
      \item Ubique en una misma recta numérica, ordenados de menor a mayor:
            \[ -\tfrac{5}{2}, \quad \sqrt{2}, \quad -1{,}5, \quad
               \tfrac{7}{4}, \quad 0{,}\overline{3}, \quad -\sqrt{4} \]
      \item ¿Verdadero o falso? Justifique cada uno con un ejemplo o un
            contraejemplo:
            \begin{itemize}\setlength{\itemsep}{-0.04cm}
              \item[a)] Todo entero es racional.
              \item[b)] Todo racional es entero.
              \item[c)] La suma de dos irracionales es siempre irracional.
              \item[d)] Entre dos racionales distintos siempre hay otro
                        racional.
            \end{itemize}
    \end{enumerate}
  \end{alertblock}
\end{frame}

\begin{frame}{Ejercicio 4 — solución}
  \scriptsize
  \begin{columns}[T,onlytextwidth]
    \begin{column}{0.42\textwidth}
      \begin{block}{Punto 1}
        \tiny
        $x = 0{,}\overline{7}$: \ $10x = 7{,}\overline{7}$, \
        $9x = 7 \Rightarrow x = \destacar{\tfrac{7}{9}}$

        \vspace{0.08cm}
        $x = 0{,}\overline{27}$: \ $100x - x = 27$, \
        $x = \tfrac{27}{99} = \destacar{\tfrac{3}{11}}$

        \vspace{0.08cm}
        $x = 1{,}\overline{3}$: \ $10x - x = 12$, \
        $x = \tfrac{12}{9} = \destacar{\tfrac{4}{3}}$
      \end{block}
    \end{column}
    \begin{column}{0.54\textwidth}
      \begin{exampleblock}{Punto 2: el orden}
        \tiny
        \[ -\tfrac{5}{2} < -\sqrt{4} < -1{,}5 < 0{,}\overline{3}
           < \sqrt{2} < \tfrac{7}{4} \]
        En decimales: $-2{,}5 < -2 < -1{,}5 < 0{,}33 < 1{,}41 < 1{,}75$.
        \vspace{0.05cm}
        {\color{uniGris} Pasar todo a decimal es la forma segura de ordenar
        cuando hay fracciones y raíces mezcladas.}
      \end{exampleblock}
    \end{column}
  \end{columns}

  \vspace{0.05cm}
  \begin{alertblock}{Punto 3}
    \tiny
    \textbf{a) Verdadero:} $n = \frac{n}{1}$. \quad
    \textbf{b) Falso:} $\frac{1}{2}$ no es entero. \quad
    \textbf{c) Falso:} $\sqrt{2} + (-\sqrt{2}) = 0$, que es racional. \quad
    \textbf{d) Verdadero:} el promedio $\frac{a+b}{2}$ siempre está entre los
    dos y sigue siendo racional. \destacar{Esa es la propiedad de densidad.}
  \end{alertblock}
\end{frame}

\begin{frame}{Errores frecuentes en esta sesión}
  \scriptsize
  \begin{center}
  \begin{tabular}{@{}p{5.0cm}p{6.8cm}@{}}
    \toprule
    \textbf{Error} & \textbf{Corrección} \\
    \midrule
    Creer que toda raíz es irracional &
      $\sqrt{16}=4$ es natural; solo lo es si el radicando no es cuadrado
      perfecto \\
    \addlinespace
    Tomar $22/7$ o $3{,}1416$ como $\pi$ &
      Son \azul{aproximaciones racionales}; $\pi$ es irracional \\
    \addlinespace
    Decir que $-5 > -2$ porque $5 > 2$ &
      En los negativos el orden se invierte: $-5 < -2$ \\
    \addlinespace
    Escribir $[-\infty,\ \infty]$ &
      El infinito nunca se cierra: $(-\infty,\ \infty)$ \\
    \addlinespace
    Calcular $|a-b|$ como $|a|-|b|$ &
      Primero se opera dentro de las barras \\
    \addlinespace
    Asumir que $-x$ es negativo &
      Es \azul{el opuesto de $x$}: si $x<0$, entonces $-x>0$ \\
    \bottomrule
  \end{tabular}
  \end{center}
\end{frame}

\begin{frame}{Cierre de la sesión}
  \begin{block}{Lo que debe quedar claro hoy}
    \scriptsize
    \begin{itemize}\setlength{\itemsep}{0.02cm}
      \item $\mathbb{N} \subset \mathbb{Z} \subset \mathbb{Q} \subset
            \mathbb{R}$, y $\mathbb{R} = \mathbb{Q} \cup \mathbb{I}$.
      \item Un número es racional si su decimal es \azul{finito o periódico};
            irracional si es infinito y no periódico.
      \item Las propiedades de $+$ y $\cdot$: clausurativa, conmutativa,
            asociativa, modulativa, invertiva y \azul{distributiva}.
      \item La distributiva leída al revés \destacar{es factorizar}: volverá
            en la unidad 3.
      \item Intervalos, desigualdades y valor absoluto como distancia al cero.
    \end{itemize}
  \end{block}

  \vspace{0.08cm}
  \begin{alertblock}{Trabajo independiente}
    \tiny
    Resolver los 20 ejercicios de la lámina siguiente. En la sesión 2
    trabajamos las \destacar{operaciones con enteros}: signos, jerarquía de
    operaciones y los errores que se arrastran todo el semestre.
  \end{alertblock}
\end{frame}

\begin{frame}{Ejercicios propuestos}
  \vspace{-0.1cm}
  \begin{beamercolorbox}[wd=\textwidth,sep=0.10cm]{block body}
    \centering\scriptsize
    \textbf{1--7:} clasificación. \quad \textbf{8--13:} propiedades y
    conversión. \quad \textbf{14--20:} orden, intervalos y valor absoluto.
  \end{beamercolorbox}
  \listatiny
  \vspace{0.15cm}
  \begin{columns}[T,onlytextwidth]
    \begin{column}{0.47\textwidth}
      \begin{enumerate}\setlength{\itemsep}{0.16cm}
        \item ¿A qué conjuntos pertenece $-\sqrt{36}$?
        \item ¿A qué conjuntos pertenece $\sqrt{10}$?
        \item ¿A qué conjuntos pertenece $\tfrac{18}{6}$?
        \item ¿$0$ es natural? ¿Es entero? ¿Es racional?
        \item Dé un irracional entre $1$ y $2$
        \item Dé un racional entre $\tfrac{1}{3}$ y $\tfrac{1}{2}$
        \item ¿Es $0{,}101001000100001\dots$ racional?
        \item Convierta $0{,}\overline{8}$ a fracción
        \item Convierta $0{,}\overline{36}$ a fracción
        \item Convierta $2{,}\overline{5}$ a fracción
      \end{enumerate}
    \end{column}
    \begin{column}{0.50\textwidth}
      \begin{enumerate}\setlength{\itemsep}{0.16cm}
        \setcounter{enumi}{10}
        \item ¿Qué propiedad es $5(x-2) = 5x-10$?
        \item ¿Qué propiedad es $(a+b)+c=a+(b+c)$?
        \item ¿Por qué $0$ no tiene inverso multiplicativo?
        \item Ordene: $-3$, $-\tfrac{7}{2}$, $\sqrt{9}$, $0$, $-\sqrt{2}$
        \item Escriba $-4 \le x < 1$ como intervalo
        \item Escriba $(2,\ \infty)$ como desigualdad
        \item Calcule $|{-12}| + |5-9|$
        \item Calcule $|3-3|$
        \item Si $x<0$, ¿cuánto vale $|x|$?
        \item ¿Es cierto que $|a+b| = |a|+|b|$ siempre?
      \end{enumerate}
    \end{column}
  \end{columns}
\end{frame}

\begin{frame}{Respuestas}{Autoevaluación de los ejercicios propuestos}
  \listatiny
  \vspace{0.15cm}
  \begin{columns}[T,onlytextwidth]
    \begin{column}{0.47\textwidth}
      \begin{enumerate}\setlength{\itemsep}{0.16cm}
        \item $-\sqrt{36}=-6$: $\mathbb{Z}$, $\mathbb{Q}$, $\mathbb{R}$
        \item Irracional: $\mathbb{I}$, $\mathbb{R}$
        \item $18/6 = 3$: $\mathbb{N}$, $\mathbb{Z}$, $\mathbb{Q}$, $\mathbb{R}$
        \item Sí, sí y sí
        \item Por ejemplo $\sqrt{2}$ o $\sqrt{3}$
        \item Por ejemplo $\tfrac{5}{12}$ (el promedio)
        \item No: infinito sin periodo, es irracional
        \item $\tfrac{8}{9}$
        \item $\tfrac{36}{99} = \tfrac{4}{11}$
        \item $\tfrac{23}{9}$
      \end{enumerate}
    \end{column}
    \begin{column}{0.50\textwidth}
      \begin{enumerate}\setlength{\itemsep}{0.16cm}
        \setcounter{enumi}{10}
        \item Distributiva
        \item Asociativa de la suma
        \item Porque no existe $y$ con $0\cdot y = 1$
        \item $-\tfrac{7}{2} < -3 < -\sqrt{2} < 0 < \sqrt{9}$
        \item $[-4,\ 1)$
        \item $x > 2$
        \item $12 + 4 = 16$
        \item $0$
        \item $|x| = -x$, que es positivo
        \item No: con $a=3$, $b=-3$ da $0 \neq 6$
      \end{enumerate}
    \end{column}
  \end{columns}
\end{frame}

\begin{frame}{Referencias}
  \footnotesize
  \begin{itemize}
    \item[{[1]}] Ramos del Olmo, Á. M. (2016). \textit{Matemáticas básicas
      para el acceso a la universidad}. Ediciones Pirámide.
      \texttt{https://elibro.net/es/ereader/uniminuto/49134}
    \item[{[2]}] González Pulido, J. W., Chacón Benavides, J. A. \&
      Fonseca Correa, L. Á. (2023). \textit{Matemática básica} (1.\textsuperscript{a} ed.).
      Universidad Mariana.
      \texttt{https://elibro.net/es/ereader/uniminuto/280647}
    \item[{[3]}] Placencia Valero, J. (2008). \textit{Compendio de matemática
      básica elemental} (1.\textsuperscript{a} ed.). Editorial Tébar Flores.
      \texttt{https://elibro.net/es/ereader/uniminuto/269758}
  \end{itemize}

  \vspace{0.35cm}
  \begin{center}
    {\Large\color{uniAzul}\textbf{¿Preguntas?}}
  \end{center}
\end{frame}

\end{document}
